\section{Projection: perspective, orthographic, isometric}
\label{sec:proj}
\here{4}
\concept{A projection maps the view volume to the canonical cube $[-1,1]^3$. Perspective shrinks distant objects
(projectors meet at the eye); parallel projections keep sizes (projectors are parallel). Multiview drawings use
orthographic views along the principal axes; an isometric view looks along $(1,1,1)$ so the three axes are
foreshortened equally.}
\begin{equation}
P_{\text{persp}}=\begin{bmatrix}\frac{f}{a}&0&0&0\\0&f&0&0\\0&0&\frac{-(F+N)}{F-N}&\frac{-2FN}{F-N}\\0&0&-1&0\end{bmatrix},\
f=\cot\tfrac{\text{fov}}{2};\qquad
P_{\text{ortho}}=\begin{bmatrix}\frac{2}{r-l}&0&0&-\frac{r+l}{r-l}\\0&\frac{2}{t-b}&0&-\frac{t+b}{t-b}\\0&0&\frac{-2}{F-N}&-\frac{F+N}{F-N}\\0&0&0&1\end{bmatrix}
\label{eq:proj}
\end{equation}
Isometric view direction: elevation $\arctan(1/\sqrt2)=35.26^\circ$, azimuth $45^\circ$; each axis is foreshortened by
$\sqrt{2/3}=0.816$.

\begin{figure}[H]
\centering
\begin{subfigure}{0.495\linewidth}\raw{scene/v10_proj_persp}{\linewidth}\caption{Perspective (fov $56^\circ$).}\end{subfigure}\hfill
\begin{subfigure}{0.495\linewidth}\raw{scene/v10_proj_iso}{\linewidth}\caption{Isometric (axonometric).}\end{subfigure}\\[1mm]
\begin{subfigure}{0.325\linewidth}\raw{scene/v10_proj_top}{\linewidth}\caption{Orthographic top (plan).}\end{subfigure}\hfill
\begin{subfigure}{0.325\linewidth}\raw{scene/v10_proj_front}{\linewidth}\caption{Orthographic front (from $+z$).}\end{subfigure}\hfill
\begin{subfigure}{0.325\linewidth}\raw{scene/v10_proj_side}{\linewidth}\caption{Orthographic side (from $+x$).}\end{subfigure}
\caption{The same scene and time under five projections (\key{Ctrl+P}). Notice in (a) that the near huts look larger
than identical far ones and parallel roads converge; in (b)--(e) equal objects stay equal anywhere in the frame.
In (c) the plan shows true distances (it is the site plan of Figure~\ref{fig:layout}); in (d) and (e) everything
is flattened onto a vertical plane, so depth is lost and only heights can be read.}
\label{fig:proj}
\end{figure}

\begin{table}[H]
\centering\small
\begin{tabularx}{\linewidth}{@{}lXXl@{}}
\toprule
Method & Visual result & Cost & When to use\\\midrule
Perspective & depth cue, convergence, realistic & same matrix cost; non-linear depth (Section~\ref{sec:projective}) & play views, video\\
Orthographic multiview & true lengths along two axes, no depth cue & same; linear depth & plans, elevations, shadow maps\\
Isometric & three faces visible, equal foreshortening & same; linear depth & layout overviews, games\\
\bottomrule
\end{tabularx}
\caption{Differences between the projections.}
\end{table}

\where{src/math/Math3D.cpp perspective / ortho; src/app/AppGame.cpp composeView (teaching projections); the sun's shadow cascades use ortho}
\critique{The screen-space effects (SSAO/GTAO, SSR, fog, TAA) reconstruct positions assuming a perspective
projection. In the orthographic views they still run, so fog and reflections in (d) and (e) are not physically
meaningful; the views are for geometry, not for lighting.}
